\section{Publish values into a flat local record}
The local record occupies 176 bytes and is eight-byte aligned. It is an application
interface, not an on-disk structure. Its integer fields and signed sample slots use
the simulator's little-endian byte order.

\tagpdfsetup{table/header-rows={1}}
\begin{center}
  \begin{tabular}{rll}
    Offset & Field             & Storage                                 \\
    \hline
    0      & Counter           & Unsigned 64 bits                        \\
    8      & Identifier length & Unsigned 32 bits                        \\
    12     & Reading count     & Unsigned 32 bits                        \\
    16     & Identifier        & Room for thirty-two bytes               \\
    48     & Readings          & Room for sixty-four signed 16-bit codes \\
  \end{tabular}
\end{center}

For the trace, the length fields become three and ten. Copy the identifier into bytes
sixteen through eighteen and the ten little-endian readings into bytes forty-eight
through sixty-seven. Leave the unused identifier and reading tails unchanged. This
makes a sentinel-filled output reveal an excessive copy or a routine that clears the
whole object instead of writing its selected fields.

No stored field points into the serialized input. The flat record can therefore be
re-encoded after the input buffer has been reused. Its used identifier bytes and sample
codes must remain intact; its unused tails have no value in the application schema and
must not influence encoded output.
